\chapter{能力 Wave 深化（Wave A--T）}
\label{ch:capability-waves}

本章归档 \versiontag{0.9.4} 之后的能力 sprint：从可证明备份（ABI v14）到备份窗口（ABI v29）。与 Part I 各章的关系：manifest v5 扩展 \cref{ch:manifest-format}；恢复/Filter 扩展 \cref{ch:restore-verify-v4}；Pipeline CDC 路由扩展 \cref{ch:streaming-cdc-v4}；CLI/daemon 子命令见 \cref{ch:daemon-cli-v4}。完整 Wave 索引见 \filepath{docs/WAVE\_ARCHIVE.md}。

\section{能力架构总览}

\begin{figure}[htbp]
\centering
\begin{tikzpicture}[
  node distance=0.55cm and 0.8cm,
  box/.style={rectangle, draw=AdrBlue!70, fill=blue!4, rounded corners=2pt,
    minimum width=2.6cm, minimum height=0.55cm, font=\scriptsize, align=center},
  side/.style={rectangle, draw=StructGreen!70, fill=green!4, rounded corners=2pt,
    minimum width=2.4cm, minimum height=0.5cm, font=\tiny, align=center},
  arr/.style={-{Stealth[length=2mm]}, thick, AdrBlue!60}
]
  \node[box] (scan) {Scan + Filter\\{\tiny scan\_entry}};
  \node[box, right=of scan] (pipe) {Pipeline CDC\\{\tiny backup\_pipeline}};
  \node[box, right=of pipe] (store) {ChunkStore\\{\tiny EbPack}};
  \node[box, right=of store] (meta) {Manifest v4/v5\\{\tiny manifest.cc}};
  \draw[arr] (scan) -- (pipe) -- (store) -- (meta);

  \node[side, below=1.1cm of scan] (pi) {path\_index.jsonl};
  \node[side, below=1.1cm of pipe] (rep) {reports/<txn>.json};
  \node[side, below=1.1cm of store] (mbi) {manifest\_browse/*.mbi};
  \node[side, below=1.1cm of meta] (job) {jobs.json / queue};

  \draw[dashed, gray] (meta.south) -- ++(0,-0.35) -| (pi.north);
  \draw[dashed, gray] (pipe.south) -- ++(0,-0.35) -| (rep.north);
  \draw[dashed, gray] (store.south) -- ++(0,-0.35) -| (mbi.north);
  \draw[dashed, gray] (meta.south) -- ++(0,-0.35) -| (job.north);

  \node[font=\scriptsize, align=center, below=0.15cm of job] {Catalog sidecar 与 Job 策略在 commit 后异步/同步落盘};
\end{tikzpicture}
\caption{能力 Wave 在备份主路径上的挂载点}
\label{fig:capability-arch}
\end{figure}

\begin{table}[htbp]
\centering
\caption{能力 Phase 与 Wave 字母对照}
\small
\begin{tabular}{@{}clll@{}}
\toprule
Phase & Wave & ABI & 主题 \\
\midrule
0 & A & — & 缺口表与路线图 \\
1--2 & — & v14--v15 & 路径索引、Merkle Diff、恢复验收 \\
3 & B, E & v16 & Manifest v5、Win 元数据保真 \\
4 & C, D & v17 & 扫描容错、备份报告、hooks \\
5 & F, G & v18 & 多 Job、GUI 作业 \\
6 & I, K, M & v19--v23 & Hybrid CDC、EBB v2、三路就地恢复 \\
7 & L & v22 & Browse sidecar、Job 队列 \\
8 & N, O, O+ & v24--v26 & 可达性、RPO、服务化、多 Profile \\
9 & P, Q & v27 & 垂直插件 \\
10 & R & v28 & 智能排除 \\
11 & T & v29 & 备份窗口 \\
12--13 & U & v30 & CompressTier + 字典 + repo-stats \\
16--19 & — & v33--v37 & 稀疏/EFS/VSS 运维/envelope/Webhook（\cref{ch:phase16-19}） \\
\bottomrule
\end{tabular}
\end{table}

%====================================================================
\section{Wave B/E：Manifest v5 与 Windows 元数据}
\label{sec:cap-manifest-v5}

\subsection{选型与触发条件}

\texttt{WriteManifestAuto(path, doc, prefer\_binary)} 在 \texttt{prefer\_binary=true} 时：
\begin{enumerate}[nosep]
  \item 若 \texttt{DocumentUsesV5(doc)} $\rightarrow$ \texttt{WriteManifestV5}（header \texttt{EBMANIFEST5}）；
  \item 否则 $\rightarrow$ \texttt{WriteManifestV4}（header \texttt{EBMANIFEST4}）。
\end{enumerate}

\texttt{DocumentUsesV5} 为 true 当任一 \texttt{ManifestFileEntry} 满足 \texttt{has\_win\_meta()}：非空 \texttt{security\_descriptor\_b64}、非零 \texttt{inode\_id}、非零 \texttt{reparse\_tag}、非空 \texttt{reparse\_target} 或 \texttt{stream\_name}。

\begin{designbox}[v4/v5 兼容策略]
读路径 \texttt{ReadManifestAuto} 对 \texttt{EBMANIFEST4} 与 \texttt{EBMANIFEST5} 共用二进制解析器；v5 body 使用 \texttt{inner\_magic = 0xEB4F0002}（v4 为 \texttt{0xEB4F0001}）。旧仓库只读兼容；新 Win 元数据备份自动升级 header。
\end{designbox}

\subsection{v5 二进制 body 扩展}

v5 在 v4 每文件 meta flags 之后追加 Win 字段（\filepath{manifest.cc}）：

\begin{longtable}{@{}llp{7.5cm}@{}}
\caption{v5 扩展 meta flags（\texttt{BuildV5MetaFlags}）} \\
\toprule
Flag & 值 & 后续字段 \\
\midrule
\texttt{kMetaWinSd} & 0x20 & \texttt{u32 sd\_len} + SD base64 bytes \\
\texttt{kMetaInode} & 0x40 & \texttt{u64 inode\_id} \\
\texttt{kMetaReparse} & 0x80 & \texttt{u32 reparse\_tag} \\
\texttt{kMetaStream} & 0x100 & \texttt{u32 stream\_len} + UTF-8 stream name \\
\texttt{kMetaReparseTarget} & 0x200 & \texttt{u32 target\_len} + UTF-8 junction target \\
\bottomrule
\end{longtable}

物理文件布局与 v4 相同：文本行 \texttt{EBMANIFEST5} + LF + binary body + 4B CRC32（LE，覆盖 body）。

\subsection{扫描 enrichment 流程}

\begin{figure}[htbp]
\centering
\begin{tikzpicture}[
  node distance=0.7cm,
  flowstep/.style={rectangle, draw, minimum width=3.2cm, minimum height=0.6cm,
    font=\scriptsize, align=center},
  arr/.style={-{Stealth}, thick}
]
  \node[flowstep, fill=blue!6] (scan) {ScanFiles 深度优先};
  \node[flowstep, fill=green!6, below=of scan] (win) {EnrichScanEntriesWinMeta\\{\tiny \_WIN32 only}};
  \node[flowstep, fill=yellow!8, below=of win] (ads) {AppendAlternateStreams\\{\tiny FindFirstStreamW}};
  \node[flowstep, fill=orange!8, below=of ads] (pipe) {Pipeline / sequential chunk};
  \draw[arr] (scan) -- (win) -- (ads) -- (pipe);
\end{tikzpicture}
\caption{Windows 扫描 enrichment（\filepath{scan\_entry.cc}）}
\label{fig:win-scan-enrich}
\end{figure}

\textbf{ADS 路径语义}：主文件 \texttt{docs/file.txt}；ADS 条目 \texttt{relative\_path = docs/file.txt:Zone.Identifier}，\texttt{stream\_name = Zone.Identifier}。恢复时 \texttt{restore\_engine} 将 \texttt{path:stream} 映射为 \texttt{CreateFileW} 写入 alternate stream。

\textbf{硬链 dedup}（Wave E）：备份阶段按 \texttt{inode\_id + stream\_name} 去重 chunk 引用；恢复阶段首路径 \texttt{CreateFile} 写内容，后续同 inode 路径 \texttt{CreateHardLinkW} 指向前者。

%====================================================================
\section{Wave C/D：尽力一致扫描与备份报告}
\label{sec:cap-scan-report}

\subsection{ScanResult 与 issue 收集}

\texttt{ScanFiles} 返回 \texttt{ScanResult\{entries[], issues[]\}}。单路径失败\textbf{不 abort} 整库：
\begin{itemize}[nosep]
  \item \texttt{unreadable}：权限/IO 错误；
  \item \texttt{depth\_exceeded}：超过 \texttt{kMaxScanDepth}；
  \item \texttt{reparse\_junction}：联接点目录不递归（WIN-03）。
\end{itemize}

Pipeline \texttt{ReaderStage} 对 issue 集合内路径 skip chunk；sequential 路径在 \texttt{ChunkPendingFiles} 读失败时同样 skip 并记入 \texttt{scan\_issues\_}。

\subsection{备份报告 JSON}

Commit 后写入 \filepath{reports/<txn>.json}（\filepath{report/backup\_report.cc}）：

\begin{longtable}{@{}lp{9cm}@{}}
\caption{备份报告主要字段} \\
\toprule
字段 & 含义 \\
\midrule
\texttt{txn\_id} & 本次事务 ID \\
\texttt{files\_processed} / \texttt{chunks\_*} & 与 \texttt{BackupStats} 对齐 \\
\texttt{scan\_issues} & 分桶计数：\texttt{unreadable}、\texttt{depth\_exceeded} 等 \\
\texttt{reparse\_junction} & 跳过的联接点目录数 \\
\texttt{hook\_failed} & pre/post hook 非零退出 \\
\texttt{plugins[]} & Wave Q：各插件 \texttt{PluginReportJson} \\
\texttt{durability\_downgraded} & Wave T：deadline 触发 Balanced \\
\texttt{window\_truncated} / \texttt{window\_end\_unix} & Wave T：超时截断 \\
\bottomrule
\end{longtable}

\subsection{Hooks}

\texttt{hook\_runner.cc} 在 backup 前后执行 shell 命令；C API \texttt{eb\_backup\_set\_backup\_hooks}；CLI \texttt{--pre-cmd}/\texttt{--post-cmd}。非零退出写入 \texttt{hook\_failed}，不 rollback 已提交 chunk（尽力一致语义）。

%====================================================================
\section{可证明备份：路径索引与 Merkle Diff（ABI v14--v15）}
\label{sec:cap-provable}

\subsection{PathIndex sidecar}

\filepath{catalog/path\_index.jsonl}：每行一条 \texttt{PathIndexRecord} JSON：
\begin{lstlisting}[caption={path\_index.jsonl 单行},label={lst:path-index-line}]
{"path":"src/main.cc","txn_id":42,"size":8192,"content_hash_hex":"...",
 "file_type":"regular","mtime_unix":1710000000}
\end{lstlisting}

\textbf{写入时机}：\texttt{AppendManifestToPathIndex} 在 manifest commit 成功后追加本次 txn 全部文件；\texttt{BuildPathIndexFromSnapshots} 可从历史 snapshot 全量 rebuild。

\textbf{查询}：\texttt{QueryPathHistoryPage(path, offset, limit)} 按 path 精确匹配分页；C API \texttt{eb\_backup\_query\_path\_history\_json}。

\subsection{Manifest Diff 与 Merkle proof}

\texttt{manifest\_diff.cc} 比较 txn A/B：
\begin{itemize}[nosep]
  \item \texttt{added} / \texttt{removed} / \texttt{modified} 文件集合；
  \item 对 \texttt{modified} 附 \texttt{merkle\_proof}（\filepath{merkle\_proof.cc} 子集证明）；
  \item 输出 \texttt{diff\_merkle\_root} 供 GUI DiffView 展示。
\end{itemize}

\subsection{恢复验收报告}

\texttt{restore\_acceptance.cc} 在 restore 完成后导出 JSON：逐文件 chunk 校验结果、Merkle 子树摘要、Win meta 应用 \texttt{issues[]}（ACL best\_effort 失败等）。

%====================================================================
\section{Wave F/G：多 Job 模型与 RunJob}
\label{sec:cap-multi-job}

\subsection{jobs.json  schema}

\begin{structbox}[BackupJob 主要字段（\filepath{job/job\_config.h}）]
\begin{itemize}[nosep]
  \item \texttt{id} / \texttt{name}：作业标识与展示名；
  \item \texttt{source\_path}：备份源根目录；
  \item \texttt{exclude\_paths[]} / \texttt{exclude\_globs[]}：Filter 合并；
  \item \texttt{plugins[]}：Wave P 插件 ID 列表；
  \item \texttt{retention\_tag} / \texttt{immutability\_days} / \texttt{worm}：Prune/WORM 策略；
  \item \texttt{window\_*}：Wave T 备份窗口（见 \cref{sec:cap-backup-window}）。
\end{itemize}
\end{structbox}

\subsection{RunJob 合并语义}

\begin{figure}[htbp]
\centering
\begin{tikzpicture}[
  node distance=0.65cm,
  box/.style={rectangle, draw, minimum width=3.4cm, minimum height=0.55cm,
    font=\scriptsize, align=center},
  arr/.style={-{Stealth}, thick}
]
  \node[box, fill=blue!6] (load) {GetJob(job\_id)};
  \node[box, fill=yellow!8, below=of load] (win) {窗外? $\rightarrow$ Conflict};
  \node[box, fill=green!6, below=of win] (merge) {合并 filter + plugins + window};
  \node[box, fill=orange!8, below=of merge] (run) {RunBackup(source, mode, opts)};
  \node[box, fill=purple!6, below=of run] (meta) {snapshot\_meta.jsonl\\job\_id / immutable\_until};
  \draw[arr] (load) -- (win) -- (merge) -- (run) -- (meta);
\end{tikzpicture}
\caption{\texttt{BackupEngine::RunJob} 控制流（\filepath{backup\_engine.cc}）}
\label{fig:run-job-flow}
\end{figure}

\texttt{snapshot\_meta.jsonl} 每 txn 一行，记录 \texttt{job\_id}、\texttt{retention\_tag}、\texttt{immutable\_until\_unix}。Prune 时 \texttt{retention\_policy.cc} 检查 WORM：\texttt{immutable\_until} 未过期则拒绝删除；immutable repo 需 ops audit key。

%====================================================================
\section{Wave I：Hybrid CDC 路由}
\label{sec:cap-hybrid-cdc}

\subsection{三路 CDC 决策}

Pipeline 对大文件（$>$ \texttt{kStreamFeedBytes}，默认 32MB）在 \textbf{full backup} 模式下按优先级路由：

\begin{table}[htbp]
\centering
\caption{CDC 路径路由（\filepath{backup\_pipeline.cc}）}
\small
\begin{tabular}{@{}llp{5.5cm}@{}}
\toprule
路径 & 启用条件 & profile 计数 \\
\midrule
Hybrid & 默认；\texttt{EBBACKUP\_CDC\_HYBRID=0} 关闭 & \texttt{hybrid\_cuts\_ns}、\texttt{hybrid\_replay\_ns} \\
Fast slice & \texttt{EBBACKUP\_CDC\_FAST\_SLICE=1} & \texttt{chunk\_ns}（整文件 cuts） \\
Stream feed & \texttt{EBBACKUP\_FORCE\_STREAM\_CDC=1} 或 fallback & \texttt{stream\_cdc\_ns} \\
\bottomrule
\end{tabular}
\end{table}

\begin{figure}[htbp]
\centering
\begin{tikzpicture}[
  decision/.style={diamond, draw, aspect=2, inner sep=1pt, font=\tiny, align=center},
  action/.style={rectangle, draw, rounded corners, minimum width=2.2cm,
    minimum height=0.5cm, font=\scriptsize},
  arr/.style={-{Stealth}, thick}
]
  \node[decision] (hybrid) {Hybrid\\enabled?};
  \node[action, below left=1cm and 0.6cm of hybrid, fill=green!8] (h) {ChunkFileStreamingHybrid};
  \node[decision, below right=1cm and 0.6cm of hybrid] (fast) {Fast slice\\enabled?};
  \node[action, below=of fast, fill=yellow!8] (f) {ChunkFileStreamingFastSlice};
  \node[action, right=1.2cm of f, fill=blue!8] (s) {ChunkFileStreaming (stream feed)};
  \draw[arr] (hybrid) -- node[left, font=\tiny]{Y} (h);
  \draw[arr] (hybrid) -- node[right, font=\tiny]{N} (fast);
  \draw[arr] (fast) -- node[left, font=\tiny]{Y} (f);
  \draw[arr] (fast) -- node[above, font=\tiny]{N} (s);
\end{tikzpicture}
\caption{单大文件 CDC 路由决策树}
\label{fig:cdc-router}
\end{figure}

Hybrid 路径：先用 \texttt{FastCdcSlice::ChunkCuts} 切分边界，再对 carry 段做 stream replay——兼顾 L5 吞吐与 chunk 边界 parity。测试断言 \texttt{stream\_cdc\_ns > 0 || hybrid\_cuts\_ns > 0}。

%====================================================================
\section{Wave M：三路就地恢复}
\label{sec:cap-three-way}

\subsection{语义：base / target / live}

就地恢复比较三份视图：
\begin{itemize}[nosep]
  \item \textbf{base}：\texttt{--base-at TXN} 或默认上一 snapshot；
  \item \textbf{target}：本次 restore 目标 txn（\texttt{--at}）；
  \item \textbf{live}：目标目录当前磁盘内容。
\end{itemize}

启用条件：\texttt{use\_three\_way=true} 且 \texttt{base\_txn\_id != target\_txn}。

\subsection{分类状态机}

\begin{longtable}{@{}llp{7cm}@{}}
\caption{三路合并 \texttt{action} 分类（\filepath{in\_place\_restore.cc}）} \\
\toprule
\texttt{action} & 条件摘要 & apply 行为 \\
\midrule
\texttt{unchanged} & live 已匹配 target & skip \\
\texttt{modify} & base$\rightarrow$target 变更，live 仍像 base & 写入 target 内容 \\
\texttt{add} & target 新增，live 无 & 创建 \\
\texttt{remove} & target 删除 & 按 orphan policy skip/delete \\
\texttt{conflict} & 路径类型冲突等 & 依 \texttt{in\_place\_conflict} 策略 \\
\texttt{both\_changed} & base 与 target 均变，live 与两者都不同 & 需用户决议（GUI bulk） \\
\bottomrule
\end{longtable}

Preview JSON 字段：\texttt{base\_action}（\texttt{same/changed}）、\texttt{live\_state}（\texttt{matches\_base/matches\_target/diverged}）、\texttt{reason}。

C API：\texttt{eb\_backup\_preview\_in\_place\_json} / \texttt{eb\_backup\_apply\_in\_place\_json}；\texttt{dry\_run=true} 仅输出计划不写盘。

%====================================================================
\section{Wave L：Manifest Browse 索引（MBI）}
\label{sec:cap-mbi}

大规模仓库 browse 不扫描完整 manifest v4/v5 body，而读 sidecar \filepath{catalog/manifest\_browse/<txn>.mbi}。

\subsection{MBI 二进制布局}

Magic \texttt{EBMBI001}（8B）；header 24B；随后 sorted records：

\begin{longtable}{@{}rrll@{}}
\caption{MBI record 序列化（\filepath{manifest\_browse\_index.cc}）} \\
\toprule
偏移 & 长度 & 字段 & 说明 \\
\midrule
0 & 4 & \texttt{path\_len u32} & UTF-8 相对路径 \\
4 & $L$ & path bytes & \\
$4+L$ & 8 & \texttt{size u64} & \\
 & 4 & \texttt{file\_type u8} + pad3 & \\
 & 8 & \texttt{mtime i64} & \\
 & 4 & \texttt{chunk\_count u32} & \\
\bottomrule
\end{longtable}

\texttt{QueryManifestBrowsePage} 按 \texttt{prefix} 二分定位 + \texttt{offset/limit} 分页；sidecar 缺失时 fallback 全量 manifest 线性扫描（\texttt{index\_source=manifest} 调试字段）。

%====================================================================
\section{Wave N/O：合规与运维可解释}
\label{sec:cap-ops}

\subsection{增量链可达性（ABI v24）}

\texttt{AnalyzeSnapshotReachability} 自 latest txn 沿 parent 链标记 reachable chunk；输出 JSON 供 \texttt{eb verify-chain} 与 GUI 徽章。不可达 chunk 不等于 orphan（可能仍被中间 snapshot 引用）——与 \texttt{GcOrphans} 语义分离。

\subsection{RPO 摘要}

\texttt{RpoSummaryJson}：最近成功 backup 时间、\texttt{stale\_backup\_alert\_days} 阈值、合规布尔；RepoHome 卡片消费。

\subsection{孤儿解释（Wave O）}

\texttt{orphan\_explain.cc}：自 \texttt{repo\_stats} + 索引枚举无 manifest 引用的 chunk，按 pack/txn 分桶；Maintenance GUI 渲染解释图；CLI \texttt{eb orphan-explain --json}。

\subsection{Ops 审计（ABI v25）}

GUI 破坏性操作（prune、GC apply 等）append-only 写入 \texttt{audit/rar.chain}；\texttt{eb audit-ops list} 分页读取。与 backup RAR 共用链格式，entry type 区分 \texttt{ops} vs \texttt{backup}。

%====================================================================
\section{Wave P/Q：垂直插件}
\label{sec:cap-plugins}

\subsection{IBackupPlugin 生命周期}

\begin{figure}[htbp]
\centering
\begin{tikzpicture}[
  node distance=0.55cm,
  phase/.style={rectangle, draw=AdrBlue!70, fill=blue!5, minimum width=2.8cm,
    minimum height=0.55cm, font=\scriptsize, align=center},
  arr/.style={-{Stealth}, thick}
]
  \node[phase] (q) {Quiesce()};
  \node[phase, below=of q] (e) {ExtraScanRoots()\\ScanHints()};
  \node[phase, below=of e] (b) {RunBackup scan/chunk};
  \node[phase, below=of b] (t) {Thaw()};
  \draw[arr] (q) -- (e) -- (b) -- (t);
  \node[font=\tiny, right=0.3cm of e, align=left] {合并额外根目录\\skip\_subtree hints};
\end{tikzpicture}
\caption{插件会话 \texttt{PluginSession} RAII 顺序}
\label{fig:plugin-lifecycle}
\end{figure}

\begin{longtable}{@{}llp{7.5cm}@{}}
\caption{内置插件} \\
\toprule
ID & 平台 & 行为 \\
\midrule
\texttt{sqlite\_checkpoint} & 全平台 & 源树 \texttt{*.db} 执行 \texttt{PRAGMA wal\_checkpoint(FULL)}；skip \texttt{-wal/-shm} \\
\texttt{registry\_hive} & Windows & \texttt{RegSaveKey} 导出至 \texttt{.ebbackup/registry/} \\
\texttt{vhdx\_scan} & Windows & \texttt{Mount-Vhd -ReadOnly} 挂载至 \texttt{.ebbackup/vhdx/<n>/} \\
\bottomrule
\end{longtable}

Wave Q：\texttt{reports/<txn>.json} 追加 \texttt{plugins[]}、\texttt{plugin\_skipped}/\texttt{plugin\_failed}；schedule JSON \texttt{plugins=} 逗号列表。

%====================================================================
\section{Wave R：智能排除建议}
\label{sec:cap-exclude}

\subsection{浅扫描算法}

\texttt{SuggestExcludeFilters(source, opts)}（\filepath{exclude\_suggestions.cc}）：
\begin{enumerate}[nosep]
  \item BFS 遍历源目录，深度 $\leq$ \texttt{max\_depth}（默认 4），访问上限 \texttt{max\_dirs\_visited}；
  \item 目录 basename 命中 catalog（\texttt{.git}、\texttt{node\_modules}、\texttt{.next}、\texttt{Thumbs.db} 等）$\rightarrow$ \texttt{exclude\_path} 建议；
  \item 文件 basename glob 命中（\texttt{*.log}）$\rightarrow$ \texttt{exclude\_glob}；
  \item \texttt{include\_ide\_dirs=false} 时跳过 \texttt{.vscode}/\texttt{.idea} 建议；
  \item 与 \texttt{existing} filter 去重；输出 \texttt{recommended} 合并集（\textbf{不自动写入} job）。
\end{enumerate}

\begin{table}[htbp]
\centering
\caption{建议条目字段}
\small
\begin{tabular}{@{}llp{6cm}@{}}
\toprule
字段 & 类型 & 说明 \\
\midrule
\texttt{apply\_as} & string & \texttt{exclude\_path} 或 \texttt{exclude\_glob} \\
\texttt{pattern} & string & 路径前缀或 glob \\
\texttt{kind} & string & \texttt{build\_artifact} / \texttt{os\_metadata} / \texttt{ide} 等 \\
\texttt{reason} & string & 人类可读解释 \\
\texttt{hit\_count} & u32 & 浅扫描命中次数 \\
\bottomrule
\end{tabular}
\end{table}

C API \texttt{eb\_backup\_suggest\_exclude\_filters\_json}；CLI \texttt{eb suggest-excludes --include-ide}。

%====================================================================
\section{Wave T：备份窗口与 deadline 自适应}
\label{sec:cap-backup-window}

\subsection{BackupWindowPolicy}

\begin{structbox}[字段语义（\filepath{job/backup\_window.h}）]
\begin{itemize}[nosep]
  \item \texttt{window\_start} / \texttt{window\_end}：本地时区 \texttt{HH:MM}（24h）；
  \item \texttt{deadline\_grace\_seconds}：距 \texttt{window\_end} 不足该秒数时触发 durability 降级（默认 300）；
  \item \texttt{durability\_adaptive}：是否启用 Strict$\rightarrow$Balanced 自适应。
\end{itemize}
\end{structbox}

\subsection{跨午夜判定}

\texttt{IsWithinBackupWindow(now, policy)}（\filepath{backup\_window.cc}）：
\begin{itemize}[nosep]
  \item 仅 \texttt{start}：\texttt{now\_min $\geq$ start\_min}；
  \item 仅 \texttt{end}：\texttt{now\_min $<$ end\_min}；
  \item \texttt{start $\leq$ end}：同日窗口 \texttt{[start, end)}；
  \item \texttt{start $>$ end}：跨午夜 \texttt{now $\geq$ start || now $<$ end}。
\end{itemize}

\begin{figure}[htbp]
\centering
\begin{tikzpicture}[x=0.12cm,y=0.8cm]
  \draw[thick,->] (0,0) -- (120,0) node[right, font=\scriptsize]{本地时间 (min)};
  \foreach \x/\lab in {0/0,30/02:00,90/06:00,120/24:00} {
    \draw (\x,0.08) -- (\x,-0.08);
    \node[below, font=\tiny] at (\x,-0.25) {\lab};
  }
  \draw[very thick, blue!70] (30,0.4) -- (90,0.4);
  \node[font=\scriptsize, blue!70] at (60,0.85) {允许备份窗口};
  \draw[fill=orange!40, opacity=0.7] (82,0.15) rectangle (90,0.65);
  \node[font=\tiny, align=center] at (86,1.1) {grace\\降级区};
  \draw[dashed, red!70] (90,0) -- (90,0.9);
  \node[font=\tiny, red!70, above] at (90,0.95) {window\_end};
\end{tikzpicture}
\caption{备份窗口与 grace 降级区（同日窗口示意）}
\label{fig:backup-window}
\end{figure}

\subsection{RunBackup 中途自适应}

\begin{enumerate}[nosep]
  \item \texttt{InitBackupWindow}：解析 policy，计算 \texttt{window\_end\_unix\_}；
  \item 每个文件 chunk 完成后 \texttt{MaybeAdaptBackupWindow}：
  \begin{itemize}[nosep]
    \item \texttt{ShouldDowngradeDurability} $\rightarrow$ \texttt{SetDurabilityMode(Balanced)}，置 \texttt{durability\_downgraded\_}；
    \item \texttt{now $\geq$ window\_end\_unix\_} $\rightarrow$ 置 \texttt{window\_truncated\_}。
  \end{itemize}
  \item 超时后 \texttt{TruncatePendingFilesAt(i)} 截断 pending 队列，commit 已完成部分；
  \item 报告写入 \texttt{durability\_downgraded}、\texttt{window\_truncated}、\texttt{window\_end\_unix}。
\end{enumerate}

\texttt{RunJob} / \texttt{RunJobQueue}：窗外返回 \texttt{Status::Conflict("outside backup window")}；queue drain \textbf{跳过}不记失败。

%====================================================================
\section{Job 队列与平台服务（Wave L/N/O+）}

\subsection{持久化队列}

\filepath{catalog/job\_queue.jsonl}：每行 \texttt{\{job\_id, mode, enqueued\_at\}}。\texttt{RunJobQueue(drain)} 顺序 pop 并 \texttt{RunJob}；daemon \texttt{queue\_drain} schedule 模式周期性 drain。

\subsection{Windows Service / systemd}

\texttt{eb service install --config PATH} 注册 SCM；\texttt{eb service run} 入口调用 \texttt{RunScheduleConfig}；循环内 \texttt{RequestDaemonStop()} + 1s 分片 sleep 支持优雅退出。Linux 模板 \filepath{deploy/ebbackup.service} + \filepath{ebbackup.timer}。

%====================================================================
\section{源码索引（能力 Wave）}

\begin{longtable}{@{}lp{9cm}@{}}
\toprule
路径 & 职责 \\
\midrule
\filepath{src/catalog/path\_index.cc} & PathIndex JSONL append/query \\
\filepath{src/catalog/manifest\_diff.cc} & Snapshot diff + Merkle \\
\filepath{src/catalog/manifest\_browse\_index.cc} & MBI sidecar 读写 \\
\filepath{src/winmeta/win\_meta.cc} & Win 元数据 capture/restore \\
\filepath{src/scan/scan\_entry.cc} & ScanResult + ADS enrichment \\
\filepath{src/scan/exclude\_suggestions.cc} & Wave R catalog + 浅扫描 \\
\filepath{src/report/backup\_report.cc} & 报告 JSON 序列化 \\
\filepath{src/common/hook\_runner.cc} & pre/post backup hooks \\
\filepath{src/job/job\_config.cc} & jobs.json CRUD \\
\filepath{src/job/backup\_window.cc} & 窗口判定 + grace \\
\filepath{src/job/job\_queue.cc} & 队列 drain \\
\filepath{src/plugin/} & 内置垂直插件 \\
\filepath{src/restore/in\_place\_restore.cc} & 三路 preview/apply \\
\filepath{src/store/orphan\_explain.cc} & 孤儿解释 JSON \\
\filepath{src/engine/backup\_engine.cc} & RunJob / window 截断 \\
\filepath{src/engine/restore\_engine.cc} & Win meta apply + hardlink \\
\filepath{src/pipeline/backup\_pipeline.cc} & Hybrid CDC 路由 \\
\filepath{src/winmeta/sparse\_file.cc} & NTFS 稀疏检测/restore（Phase 16） \\
\filepath{src/winmeta/efs\_key.cc} & EFS Read/WriteEncryptedFileRaw（Phase 19） \\
\filepath{src/winmeta/vss\_shadow\_storage.cc} & WMI 影子存储查询（Phase 18） \\
\filepath{src/crypto/envelope.cc} & crypto.envelope.json（Phase 17） \\
\filepath{src/engine/manifest.cc} & kMetaSparse / kMetaEfs 读写 \\
\filepath{cli/eb\_service\_win.cc} & Windows SCM 集成 \\
\bottomrule
\end{longtable}
